% Generated draft: requires mathematical and source review.
\documentclass[11pt]{article}
\usepackage[margin=25mm]{geometry}
\usepackage{fontspec}
\setmainfont{Latin Modern Roman}
\usepackage{amsmath,amssymb,mathtools,unicode-math}
\setmathfont{Latin Modern Math}
\usepackage{enumitem,needspace,xurl,hyperref}
\setlength{\parindent}{0pt}
\setlength{\parskip}{5pt}
\setlength{\emergencystretch}{3em}
\newcommand{\N}{\mathbb{N}}
\newcommand{\Z}{\mathbb{Z}}
\newcommand{\Q}{\mathbb{Q}}
\newcommand{\R}{\mathbb{R}}
\newcommand{\C}{\mathbb{C}}
\newcommand{\F}{\mathbb{F}}
\newcommand{\A}{\mathbb{A}}
\newcommand{\PP}{\mathbb{P}}
\newcommand{\E}{\mathbb{E}}
\newcommand{\eps}{\varepsilon}
\newcommand{\dd}{\,\mathrm{d}}
\begin{document}
\section*{UnsolvedMath ID: 701596}
\textbf{Source problem code:} catalogue entry
\subsection{Definitions and mathematical statement}
A full net (also called a carpet) of order \(n\) over a field \(K\) is a family \(\sigma=(\sigma_{ij})_{1\leq i,j\leq n}\) of additive subgroups of \(K\) such that
\[
\{ab:a\in\sigma_{ir},\ b\in\sigma_{rj}\}\subseteq\sigma_{ij}
\qquad(1\leq i,r,j\leq n).
\]
The source calls a net irreducible if all its constituent subgroups are nontrivial. It calls a net that ``does not contain the diagonal'' an elementary net, but does not specify precisely how diagonal entries and the multiplication conditions are handled in this definition.

For an integer \(n\geq1\), let \(I_n\) be the identity matrix and let \(e_{ij}\) be the matrix unit. For \(i\ne j\) and \(a\in K\), define the elementary transvection
\[
t_{ij}(a)=I_n+a e_{ij}.
\]
Given the off-diagonal subgroups of \(\sigma\), define its elementary net subgroup by
\[
E(\sigma)=\bigl\langle t_{ij}(a):1\leq i,j\leq n,\ i\ne j,\ a\in\sigma_{ij}\bigr\rangle
\leq\operatorname{SL}_n(K).
\]
The requirement that this group contain no additional elementary transvections is expressed by
\[
\forall i\ne j\ \forall a\in K,\qquad
 t_{ij}(a)\in E(\sigma)\ \Longleftrightarrow\ a\in\sigma_{ij}.
\tag{*}
\]
An elementary net satisfying \((*)\) is called closed.

\textbf{Question.} Let \(R\) be a principal ideal domain with identity, let \(k=\operatorname{Frac}(R)\), and let \(K/k\) be an algebraic field extension. Regard \(R\) as a subring of \(K\) through \(R\subseteq k\subseteq K\). Let \(\sigma\) be an irreducible elementary net over \(K\), and assume that every constituent subgroup \(\sigma_{ij}\) is an \(R\)-submodule of \(K\), meaning that
\[
r a\in\sigma_{ij}\qquad(r\in R,\ a\in\sigma_{ij}).
\]
Must \(\sigma\) be closed, that is, must \((*)\) hold?

\textbf{Unresolved qualification.} The closure target \((*)\) is explicit, but the supplied definition does not unambiguously determine the class of elementary nets quantified over. In particular, it does not say whether diagonal subgroups are omitted or assigned values, which multiplication conditions survive when diagonal entries are omitted, or whether irreducibility then concerns only off-diagonal entries. No diagonal convention is imposed here. The source also supplies no restriction on the order \(n\) beyond its use as the size of the net.
\subsection{Short English statement}
Over an algebraic extension of the fraction field of a principal ideal domain, must every irreducible elementary net whose entries are modules over that domain generate no elementary transvections beyond its prescribed entries? The supplied diagonal convention is incomplete, so the admissible class of nets requires clarification.
\subsection{Sources}
\begin{itemize}
\item Local UnsolvedMath registry, permanent ID \(701596\), exact source problem code \texttt{catalogue entry}, title \emph{Kourovka Notebook Problem 19.48}; supplied complete source record, upstream revision \texttt{d371ee917ac9647fbb06cd6fc529deec6f3dda82}.
\item The supplied attribution identifies \emph{The Kourovka Notebook (Unsolved Problems in Group Theory)}, 19th issue (2018), Problem 19.48, and the 21st issue edition (2026), \texttt{arXiv:1401.0300v48}. These bibliographic details are supplied source data and have not been independently checked.
\item The proposer field is truncated as supplied: \texttt{V. A. Ko\textbackslash u\{\textbackslash i}. The background separately contains the disconnected fragment \texttt{baev\}}. A complete proposer name cannot reliably be recovered from this record.
\item The registry labels the problem \texttt{open} and gives the proposed year as 2018. This is a supplied status label, not an independently verified assessment of the literature. No literature-verification date or separate literature-source links are supplied.
\end{itemize}
\end{document}
